\documentclass[11pt,a4paper]{article}
\usepackage{iftex}
\ifPDFTeX
\usepackage[utf8]{inputenc}
\usepackage[T1]{fontenc}
\usepackage{lmodern}
\else
\usepackage{fontspec}
\setmainfont{lmroman10-regular.otf}[BoldFont=lmroman10-bold.otf,ItalicFont=lmroman10-italic.otf,BoldItalicFont=lmroman10-bolditalic.otf]
\setmonofont{lmmono10-regular.otf}
\fi
\usepackage[ngerman]{babel}
\usepackage[a4paper,margin=24mm]{geometry}
\usepackage{microtype}
\usepackage{amsmath,amssymb}
\usepackage{booktabs,longtable,tabularx,array}
\usepackage{enumitem}
\usepackage{xcolor}
\usepackage{xurl}
\usepackage{hyperref}
\definecolor{navy}{HTML}{12334A}
\definecolor{teal}{HTML}{146B75}
\hypersetup{colorlinks=true,linkcolor=navy,citecolor=teal,urlcolor=teal,pdftitle={Was ESG-Scores messen und wie sich Greenwashing prüfen lässt},pdfauthor={ETHack 2026 -- Forschungs- und Methodenbericht}}
\urlstyle{same}
\newcolumntype{P}[1]{>{\raggedright\arraybackslash}p{#1}}
\newcommand{\coeq}{\ensuremath{\mathrm{CO}_{2}\mathrm{e}}}
\newcommand{\tonnen}{\ensuremath{\mathrm{t\,CO}_{2}\mathrm{e}}}
\newcommand{\NA}{\textnormal{nicht bestimmbar}}
\newcommand{\doi}[1]{\href{https://doi.org/#1}{\nolinkurl{doi:#1}}}
\setlength{\parindent}{0pt}
\setlength{\parskip}{5pt}
\setlength{\tabcolsep}{4pt}
\setlength{\emergencystretch}{2em}
\setlist{itemsep=3pt,topsep=4pt,leftmargin=*}
\renewcommand{\arraystretch}{1.17}
\setcounter{tocdepth}{2}
\title{\textcolor{navy}{Was ESG-Scores messen\vspace{2mm}\\und wie sich Greenwashing prüfen lässt}\\[5mm]\large Aussagekräftige Kennzahlen, Evidenzgrenzen und ein Forschungsdesign für den S\&P~500}
\author{Forschungs- und Methodenbericht für ETHack 2026}
\date{Recherchestand: 12. September 2026}

\begin{document}
\maketitle
\begin{abstract}
Ein ESG-Rating ist weder ein physikalisches Nachhaltigkeitsmaß noch ein Nachweis ehrlicher Unternehmenskommunikation. Viele etablierte Ratings bewerten vor allem die finanziellen Risiken von Nachhaltigkeitsthemen für Unternehmen. Die gesellschaftlichen Schäden ihrer Geschäftstätigkeit sind ein anderes Messobjekt. Dieser Bericht untersucht die empirische Aussagekraft von ESG-Ratings, entwickelt ein dimensionskonsistentes Kennzahlensystem und operationalisiert die Prüfung von Klimaaussagen. Im Zentrum stehen Bruttoemissionen mit konsistenten Bilanzgrenzen, physische Intensitäten, kumulative Emissionspfade, ausgeführte Investitionen sowie die Trennung von Strombilanzierung, Kompensation und zusätzlicher Klimawirkung. Greenwashing wird nicht aus einem niedrigen Score oder fehlenden Angaben abgeleitet, sondern anhand spezifischer Behauptungen und abgestufter Belege untersucht. Vorgeschlagen wird ein \emph{Claim--Evidence--Outcome Audit}: ein prüfbares System, das Unternehmen vergleicht, dokumentierte Inkonsistenzen identifiziert und bei unzureichender Evidenz ausdrücklich auf ein Urteil verzichtet. Es handelt sich um eine Literatur- und Methodensynthese, nicht um eine bereits durchgeführte Bewertung aller S\&P-500-Unternehmen.
\end{abstract}
\vfill
\textbf{Leseschlüssel.} Literaturbefunde, methodische Vorschläge und synthetische Rechenbeispiele werden getrennt gekennzeichnet. Der Bericht enthält keine Feststellung, dass ein bestimmtes Unternehmen täuscht, keine berechneten Portfoliogewichte und keine behauptete empirische Modellgüte.
\newpage
\tableofcontents
\newpage

\section{Management Summary: die entscheidenden Antworten}
\label{sec:summary}
\begin{enumerate}
\item \textbf{Wie aussagekräftig ist ein ESG-Score?} Nur relativ zu seiner definierten Zielgröße. Ein Rating zur finanziellen Resilienz kann nützlich sein, ohne niedrige Emissionen, geringe lokale Schäden oder zusätzliche Klimawirkung nachzuweisen. Anbieterübereinstimmung ist eine Frage der Reliabilität; ökologische Richtigkeit eine andere Frage der Validität.
\item \textbf{Welche Größen sind wirklich aussagekräftig?} Ergebnisnahe Größen mit expliziter Systemgrenze: absolute Bruttoemissionen, physische Emissionsintensität, reale Fortschritte auf vergleichbarer Anlagenbasis, standortbezogene Umweltbelastung und dokumentierte schwere soziale Schäden. Keine dieser Größen ist ohne Kontext ein universeller Nachhaltigkeitsmaßstab.
\item \textbf{Welche Größen sind nur Indizien?} Ziele, Policies, ESG-verknüpfte Vergütung, Berichtsumfang, Zertifikate und angekündigte Investitionen. Sie können künftige Ergebnisse erklären, dürfen diese aber nicht ersetzen. Ihre Prognosekraft muss getestet werden.
\item \textbf{Wann liegt belastbare Greenwashing-Evidenz vor?} Wenn eine konkret abgegrenzte, wesentliche Aussage durch passende Belege nicht gestützt oder widerlegt wird. Eine Datenlücke ist zunächst eine Datenlücke; ein Zukunftsziel ist keine Behauptung bereits erreichter Reduktionen. Vorsatz und rechtliche Irreführung folgen nicht automatisch aus einer statistischen Abweichung.
\item \textbf{Was wäre ein substanzieller Hackathon-Beitrag?} Nicht ein weiterer gewichteter ESG-Gesamtscore, sondern eine reproduzierbare Prüfung, welche behaupteten Verbesserungen nach Bereinigung um Preiseffekte, Bilanzgrenzen und Bilanzierungsinstrumente noch bestehen bleiben. Dazu kommen Gegenbelege, Unsicherheitsbereiche und eine explizite Klasse \glqq nicht entscheidbar\grqq.
\item \textbf{Was bedeutet dies für einen Net-Zero-Fonds?} Finanzierte Emissionen sind eine Zurechnungsgröße. Für tatsächlichen Investoreneinfluss braucht es zusätzlich einen plausiblen Wirkungsmechanismus, etwa wirksames Engagement oder zusätzliche Finanzierung. Portfolio-Umschichtung allein belegt keine vermiedene Tonne.
\end{enumerate}

\section{Fragestellung, Begriffe und Evidenzstandard}
\subsection{Drei Messobjekte, die nicht verwechselt werden dürfen}
Die ETHack-Aufgabe verlangt einen datenbasierten Nachhaltigkeitsvergleich von S\&P-500-Unternehmen und optional eine begründete Allokation eines Fonds von 1~Mrd.~USD unter einem beschleunigten Net-Zero-Szenario \cite{deck}. Dafür müssen drei unterschiedliche Fragen beantwortet werden:
\begin{description}
\item[Finanzielle Materialität:] Wie wirken Umwelt- und Sozialthemen auf Cashflows, Kapitalkosten und Unternehmenswert? Beispiel: Ein Wasserengpass unterbricht die Produktion.
\item[Wirkungsmaterialität:] Wie beeinflusst das Unternehmen Menschen und Umwelt? Beispiel: Wasserverbrauch verschärft einen lokalen Nutzungskonflikt, auch wenn das Unternehmen dafür nichts bezahlt.
\item[Zusätzliche Wirkung einer Intervention:] Was verändert eine konkrete Maßnahme gegenüber einer plausiblen Welt ohne diese Maßnahme? Beispiel: Ermöglicht ein Finanzierungsvertrag eine ansonsten nicht gebaute Anlage?
\end{description}
Ein Unternehmen kann finanziell resilient und gesellschaftlich schädlich sein. Ein Unternehmen mit heute hohen Emissionen kann reale Reduktionen erzielen. Umgekehrt kann ein emissionsarmes Geschäftsmodell ohne weitere Verbesserungen bestehen. Diese Aussagen sind logisch miteinander vereinbar.

\textbf{Arbeitsdefinition.} Nachhaltigkeitsleistung bezeichnet hier das Niveau und die Veränderung wesentlicher Umwelt- und Sozialwirkungen unter definierten Systemgrenzen. Klimaleistung wird am weitesten operationalisiert. Wasser, Schadstoffe, Arbeitssicherheit und Menschenrechte werden als eigene Dimensionen ergänzt; sie werden nicht durch einen Klimascore ersetzt.

\subsection{Art und Grenzen dieses Berichts}
Grundlage ist eine gezielte, nicht als vollständig beanspruchte Recherche von Fachartikeln, institutionellen Working Papers, Standards und offiziellen Datenbeschreibungen. Bevorzugt wurden Primärquellen mit expliziten Messmethoden und identifizierbarer Stichprobe. Anbieterbeschreibungen belegen, \emph{was ein Produkt zu messen beansprucht}; sie sind keine unabhängige Validierung seiner Güte. Working Papers werden als solche bezeichnet.

Die Literatur wird weder neu metaanalysiert noch anhand eines vollständigen S\&P-500-Panels repliziert. Quantitative Literaturbefunde gelten für die jeweils untersuchten Daten, Zeiträume und Methoden. Die eigenen Formeln und Entscheidungsregeln sind \textbf{zu validierende Designvorschläge}, keine bereits wissenschaftlich bestätigten Ratings. Alle folgenden synthetischen Beispiele dienen ausschließlich der Prüfung logischer Eigenschaften.

\section{ESG-Scores: Informationsgehalt und systematische Grenzen}
\label{sec:esg}
\subsection{Was etablierte Anbieter tatsächlich bewerten}
MSCI beschreibt seine Ratings als branchenrelative Bewertung der Resilienz gegenüber finanziell relevanten Nachhaltigkeitsrisiken und -chancen; die Buchstabenklassen reichen von AAA bis CCC \cite{msci}. Sustainalytics beschreibt den Gegenstand als ungemanagtes, durch ESG-Faktoren verursachtes wirtschaftliches Risiko \cite{sustainalytics}. Die Produkte haben damit verwandte, aber nicht identische Zielgrößen und Skalen. Bei einem Risikoscore kann ein niedriger Wert günstiger sein, während bei einem Leistungsscore ein hoher Wert günstiger ist.

\textbf{Folgerung:} Ein gutes Risikorating bedeutet nicht notwendigerweise geringe absolute Umweltbelastung. Diese Nichtgleichheit ist zunächst kein Versagen des Ratings, sondern möglicherweise eine Fehlverwendung durch seinen Nutzer. Auch ein branchenführendes emissionsintensives Unternehmen kann mehr Schäden verursachen als ein mittelmäßig bewertetes Unternehmen einer anderen Branche.

\subsection{Warum ein ESG-Gesamtscore kein Naturmaß ist}
Ein vereinfachtes Messmodell lautet
\begin{equation}
R_{i,a,t}=f_{a,v}\!\left(X_{i,t},D_{i,t},P_{i,t};J_{a,v},w_{a,v}\right),
\end{equation}
wobei $X$ beobachtete Ergebnisse, $D$ Offenlegung, $P$ Managementprozesse, $J$ die ausgewählten Themen, $w$ normative Gewichte und $v$ die Methodenversion bezeichnet. Anbieter $a$ können zusätzlich unterschiedliche Vergleichsgruppen und Datenquellen verwenden.

Viele Gesamtscores sind eher \emph{formative Indizes}: Unterschiedliche Merkmale definieren zusammen den Index. Sie sind nicht notwendigerweise Messungen einer einzigen latenten Eigenschaft, die sämtliche Indikatoren gemeinsam verursacht. Geringe Korrelation zwischen Arbeitssicherheit und Emissionen ist deshalb nicht automatisch ein Messfehler. Umgekehrt macht hohe interne Korrelation einen Index noch nicht ökologisch gültig.

Aggregation setzt außerdem Kompensierbarkeit voraus. Eine gewichtete Summe kann bessere Governance gegen höhere Emissionen verrechnen. Ob ein solcher Tausch gesellschaftlich akzeptabel ist, entscheidet keine Statistik. Für bestimmte schwere Menschenrechtsverletzungen oder bindende Grenzwerte sind separat ausgewiesene Ausschluss- beziehungsweise Prüfregeln sachgerechter als zusätzliche Minuspunkte.

Buchstabenklassen sind Ordnungen, keine physikalischen Abstände: AAA ist nicht \glqq doppelt so nachhaltig\grqq{} wie eine niedrigere Klasse. Auch numerische Anbieterscores dürfen nur entsprechend ihrer dokumentierten Skala interpretiert werden. Ein Branchenperzentil kann steigen, weil andere Firmen schlechter werden; es beweist keinen absoluten Fortschritt.

\subsection{Was die empirische Forschung zeigt -- und was nicht}
\small
\begin{longtable}{@{}P{0.22\linewidth}P{0.39\linewidth}P{0.33\linewidth}@{}}
\toprule
\textbf{Quelle} & \textbf{Befund} & \textbf{Zulässige Interpretation und Grenze} \\
\midrule
\endfirsthead
\toprule
\textbf{Quelle} & \textbf{Befund} & \textbf{Zulässige Interpretation und Grenze} \\
\midrule
\endhead
Berg et al.\ (2022) \cite{berg2022} & Bei sechs Anbietern entfallen 56\,\% der Ratingdivergenz auf Messung, 38\,\% auf Themenabdeckung und 6\,\% auf Gewichtung. & Dekomposition von \emph{Uneinigkeit}, nicht von Fehlern gegenüber einer beobachteten Nachhaltigkeitswahrheit. Gewichtsoptimierung allein löst das Problem nicht. \\
Christensen et al.\ (2022) \cite{christensen2022} & Mehr ESG-Offenlegung geht mit größerer Bewertungsdivergenz einher; die Arbeit untersucht auch Offenlegungsschocks. & Mehr Information erzeugt ohne gemeinsame Interpretation nicht zwingend mehr Einigkeit. Kein Argument dafür, weniger offenzulegen. \\
Drempetic et al.\ (2020) \cite{drempetic2020} & Unternehmensgröße, Ressourcen, Datenverfügbarkeit und ASSET4-ESG-Scores hängen positiv zusammen. & Berichtskapazität kann das Ranking beeinflussen. Der Befund beweist nicht, dass jeder Größeneffekt ausschließlich Messbias ist. \\
Chatterji et al.\ (2009) \cite{chatterji2009} & KLD-Umweltbedenken enthalten Information über frühere und teilweise spätere Verschmutzung und Verstöße; Umweltstärken prognostizieren diese Ergebnisse nicht entsprechend. & Negative Ereignisse und positive Policies sind nicht symmetrische Signale. Historischer US-/KLD-Befund, kein Urteil über jedes heutige Rating. \\
Raghunandan und Rajgopal (2022) \cite{raghunandan2022} & US-ESG-Fonds von 2010--2018 halten in ihren Vergleichen Firmen mit höheren ESG-Scores, aber schlechteren bestimmten Umwelt-/Arbeitsrechtsbilanzen und höherer Umsatz-Emissionsintensität. & Ein ESG-Fondslabel garantiert keine bessere Stakeholder-Bilanz. Ergebnis hängt von Fondsauswahl, Vergleichsgruppe und Beobachtungszeitraum ab. \\
Treepongkaruna et al.\ (2024) \cite{treepongkaruna2024} & Für US-Unternehmen 2005--2018 bedeuten hohe Refinitiv-ESG-/E-Scores in der Studie nicht niedrigere CO$_2$-Emissionen. & Rating und Emissionsleistung können auseinanderfallen. Granger-Vorhersagetests belegen weder Täuschungsabsicht noch einen kausalen Effekt des Ratings. \\
Khan et al.\ (2016) \cite{khan2016} & Gute Bewertungen finanziell materieller Nachhaltigkeitsthemen hängen mit besserer späterer Finanzperformance zusammen. & Positive Evidenz für bestimmte finanzielle Anwendungen; kein Nachweis geringerer Externalitäten und keine garantierte zukünftige Überrendite. \\
Berg et al.\ (Working Paper, 2021) \cite{berg2021} & Rückwirkende Änderungen historischer Refinitiv-Daten verändern die gemessenen Zusammenhänge zwischen ESG und Renditen. & Datenversionen und damalige Verfügbarkeit sind für Backtests wesentlich. Eine nachträgliche Revision kann auch eine legitime Korrektur sein. \\
\bottomrule
\end{longtable}
\normalsize

\subsection{Aussagekraft nach Verwendungszweck}
\textbf{Als Risiko-Screening:} potenziell nützlich, wenn Zielgröße, Branche und Methodik passen. \textbf{Als universelle Rangliste gesellschaftlicher Güte:} nicht hinreichend validiert. \textbf{Als Greenwashing-Detektor:} ungeeignet ohne konkrete Aussageprüfung. \textbf{Als Ersatz für Emissionsdaten:} ungeeignet. \textbf{Als Forschungsvariable:} verwendbar, sofern Datenstand, Anbieterabhängigkeit, Messfehler und Endogenität berücksichtigt werden.

Edmans argumentiert entsprechend, ESG als Teil normaler langfristiger Unternehmensanalyse ernst zu nehmen, statt jede ESG-bezogene Größe als besondere Qualitätsgarantie zu behandeln \cite{edmans2023}. Die sachgerechte Antwort lautet daher weder \glqq ESG ist bedeutungslos\grqq{} noch \glqq ESG misst Nachhaltigkeit\grqq, sondern: \textbf{Welches Konstrukt wird für welchen Zweck wie gut gemessen?}

\section{Welche messbaren Größen tragen tatsächlich Information?}
\label{sec:metrics}
\subsection{Auswahlprinzip: Nähe zum Ergebnis und dokumentierte Grenzen}
Ein Indikator ist besonders nützlich, wenn er (i) ein wesentliches Ergebnis abbildet, (ii) eine nachvollziehbare Einheit und Grenze hat, (iii) zeitlich vergleichbar ist, (iv) nicht leicht durch bloßes Reporting verbessert werden kann und (v) unabhängig überprüfbar ist. Auch physische Daten können geschätzt, unvollständig oder strategisch abgegrenzt sein. \glqq Quantitativ\grqq{} ist deshalb kein Synonym für \glqq wahr\grqq.

\subsection{Bruttoemissionen: ein Vektor statt einer vermeintlich vollständigen Zahl}
Für Unternehmen $i$ und Berichtsjahr $t$ ist zunächst der Vektor
\begin{equation}
\mathbf E_{it}=\left(E^{(1)}_{it},E^{(2,LB)}_{it},E^{(2,MB)}_{it},E^{(3,1)}_{it},\ldots,E^{(3,15)}_{it}\right)
\end{equation}
zu erfassen. Alle Komponenten bezeichnen \emph{jährliche Emissionsmengen} in \tonnen. LB und MB stehen für location-based und market-based. Sie sind alternative Scope-2-Bilanzierungen und werden \textbf{niemals addiert}. Die Scope-3-Kategorien werden mit ihrer Abdeckung und Berechnungsmethode gespeichert \cite{ghgcorporate,ghgscope2,ghgscope3}.

Eine Gesamtsumme ist nur bei kompatiblen Grenzen und hinreichender Abdeckung sinnvoll. Eine partielle Summe darf nicht als vollständiger Fußabdruck erscheinen. Bilanzierte biogene Emissionen, Entnahmen, gespeicherter Kohlenstoff und Kompensationszertifikate benötigen jeweils separate Behandlung nach dem verwendeten Standard; sie werden nicht stillschweigend in Bruttowerte eingerechnet.

Busch et al.\ zeigen erhebliche Datenkonsistenzprobleme, insbesondere bei Scope~3 \cite{busch2022}. Klaaßen und Stoll finden in ihrer Fallstudie von 56 Technologieunternehmen große Lücken zwischen berichteten und harmonisierten Fußabdrücken \cite{klaassen2021}. Der konkrete Befund ist keine pauschale Schätzregel für alle Firmen. Insbesondere darf nicht jedes berichtete Scope~3 mechanisch verdoppelt werden.

\textbf{Aggregationsgrenze:} Die Emission eines Stromerzeugers kann gleichzeitig Scope~2 eines Kunden und Scope~3 eines weiteren Unternehmens sein. Das ist innerhalb unterschiedlicher Unternehmensinventare erwartbar. Die Summe aller S\&P-500-Inventare ist aber dadurch keine überschneidungsfreie globale Emissionsmenge.

\subsection{Physische Intensität und absolute Veränderung gemeinsam betrachten}
Bei gleichem sachlichem Umfang und positivem Produktionsvolumen $Q_{it}$ gilt
\begin{equation}
I^{\mathrm{phys}}_{it}=\frac{E_{it}}{Q_{it}}.
\end{equation}
Geeignete Einheiten sind beispielsweise t~\coeq/MWh für Stromerzeugung, t~\coeq/t Rohstahl oder kg~\coeq/Passagier-km. Unterschiedliche Produkte und Qualitäten sind nicht beliebig austauschbar; Konglomerate benötigen Segmentdaten. Emissionen des gesamten Konzerns dürfen nicht durch die Produktion eines einzelnen Segments dividiert werden.

Für positive Werte und konsistente Grenzen gilt die exakte Identität
\begin{equation}
\Delta\ln E=\Delta\ln Q+\Delta\ln I^{\mathrm{phys}}.
\label{eq:decomposition}
\end{equation}
Sie trennt Aktivitätsänderung und Intensitätsänderung. Bei steigender Produktion kann die Intensität sinken, während absolute Emissionen steigen. Umgekehrt ist eine absolute Reduktion durch einen konjunkturellen Produktionseinbruch real, aber nicht automatisch eine strukturelle Verbesserung der Technologie.

Eine relative Veränderung ist nur für $E_{i0}>0$ definiert:
\begin{equation}
r_i=1-\frac{E_{it}}{E_{i0}}.
\end{equation}
Zusätzlich wird stets die absolute Differenz $E_{i0}-E_{it}$ ausgewiesen. Bei Basiswert null, unklarer Bilanzgrenze oder fehlendem Wert wird keine Prozentzahl erzeugt. Hohe prozentuale Reduktionen bei winzigen Ausgangswerten sind nicht automatisch größere gesellschaftliche Beiträge.

\subsection{Warum CO\texorpdfstring{$_2$}{2} pro Umsatz leicht täuscht}
Für Umsatz $R_{it}>0$ ist $I^{\mathrm{rev}}_{it}=E_{it}/R_{it}$ eine ökonomische Intensität, keine technische Effizienz. Es gilt
\begin{equation}
\Delta\ln I^{\mathrm{rev}}=\Delta\ln E-\Delta\ln R.
\end{equation}
Preise, Wechselkurse, Produktmix und Inflation verändern den Nenner. Eine Deflationierung kann allgemeine Preiseffekte reduzieren, beseitigt aber keine unternehmensspezifischen Preis- und Mixeffekte. Umsatzintensität bleibt ein ergänzender Vergleich, insbesondere bei fehlenden physischen Produktionsdaten. Für nichtpositive oder fehlende Nenner lautet das Ergebnis \NA; EBITDA ist wegen negativer Werte und Branchenunterschieden kein universell besserer Nenner.

\subsection{Bilanzgrenzen: Dekarbonisierung oder Verkauf des Problems?}
Die Prüfung benötigt eine \emph{Überleitungsrechnung}, keine universelle Korrekturformel mit pauschalem Addieren von Akquisitionen und Subtrahieren von Verkäufen. Zwei Perspektiven sind getrennt relevant:
\begin{enumerate}
\item \textbf{Aktueller Unternehmensumfang:} historisches Basisjahr auf den heutigen Konsolidierungskreis zurückrechnen, soweit Daten vorliegen.
\item \textbf{Konstante Anlagenkohorte:} dieselben Anlagen vor und nach Eigentümerwechsel verfolgen. Für eine feste Menge $\mathcal A_0$ ist $E^{\mathrm{assets}}_t=\sum_{a\in\mathcal A_0}E_{a,t}$ unabhängig davon zu beobachten, wem diese Anlagen inzwischen gehören.
\end{enumerate}
Die erste Perspektive unterstützt einen fairen Unternehmensvergleich; die zweite untersucht, ob Verschmutzung tatsächlich verschwindet oder nur den Eigentümer wechselt. Produktionsverlagerungen und Outsourcing können darüber hinaus eine Lieferkettenprüfung erfordern. Ohne Überleitung wird ein Trend als \glqq nicht vergleichbar\grqq{} markiert, nicht automatisch als Greenwashing.

Duchin, Gao und Xu untersuchen Verkäufe umweltbelastender US-Anlagen. Ihre Befunde zeigen, dass Verkäufer bessere ESG-Bewertungen erhalten können, ohne dass die Verschmutzung der veräußerten Anlagen entsprechend sinkt \cite{duchin2025}. Das liefert ein konkretes Forschungsdesign: \textbf{Unternehmensrating und anlagenbezogene Ergebnisse parallel verfolgen.}

\subsection{Scope 2: Bilanzierungsdifferenz ist kein kausaler Effekt}
Die Differenz $E^{(2,LB)}-E^{(2,MB)}$ zeigt den Einfluss unterschiedlicher Bilanzierungsansätze, nicht direkt die Menge vermiedener Emissionen. Eine große Differenz ist ein Prüfhinweis, kein automatischer Täuschungsnachweis. Location-based verwendet durchschnittliche Netzemissionsfaktoren und ist ebenfalls kein unmittelbares Maß marginaler oder kausal vermiedener Emissionen.

Bjørn et al.\ zeigen für die untersuchten Unternehmen 2015--2019, dass die Herausrechnung REC-basierter Reduktionsansprüche die Bewertung der Zielkompatibilität verändert. Die im Artikel genannte Größenordnung von 42\,\% betrifft zukünftige zugesagte Scope-2-Reduktionen \emph{unter Fortsetzung des beobachteten Trends}, nicht alle Stromzertifikate weltweit \cite{bjorn2022}.

Ein langfristiger Beschaffungsvertrag, ein ungebündeltes Zertifikat und der Bau zusätzlicher Erzeugung sind unterschiedliche Interventionen. Für zusätzliche Wirkung sind unter anderem Finanzierungseinfluss, Ort, Zeitpunkt, Erzeugungsprofil und ein plausibles Ohne-Maßnahme-Szenario relevant. \glqq Erneuerbarer Strom eingekauft\grqq{} und \glqq zusätzliche fossile Erzeugung verdrängt\grqq{} sind verschiedene Behauptungen.

\subsection{Zielabdeckung, kumulative Pfade und tatsächliche Umsetzung}
Ein Ziel muss mindestens Basisjahr, Zieljahr, erfasste Scopes, absolute oder intensitätsbezogene Form und die abgedeckten Aktivitäten angeben. Bei bekanntem vollständigem Inventar wäre der emissionsgewichtete Abdeckungsanteil
\begin{equation}
c_i=\frac{\sum_{k\in\mathcal K_i^{\mathrm{Ziel}}}E_{ik}}{\sum_{k\in\mathcal K_i^{\mathrm{wesentlich}}}E_{ik}}.
\end{equation}
Ist der Nenner unvollständig, ist $c_i$ nicht zuverlässig bestimmbar. Die Quote \glqq 10 von 15 Kategorien berichtet\grqq{} darf alternativ als \emph{Kategorienabdeckung} erscheinen, nicht als \glqq zwei Drittel der Emissionen abgedeckt\grqq.

Net Zero betrifft nicht nur einen Endpunkt. Für jährliche Mengen, ein Szenario $s$ und den Horizont $T$ seien
\begin{align}
C_i^{(s)}&=\sum_{\tau=t_0}^{T}E_{i\tau}^{(s)},\\
B_i^{(s)}&=\sum_{\tau=t_0}^{T}Q_{i\tau}^{(s)}I_{\mathrm{Sektor},\tau}^{(s)},\\
A_i^{(s)}&=\frac{C_i^{(s)}}{B_i^{(s)}}\quad\text{für }B_i^{(s)}>0.
\label{eq:budget}
\end{align}
$C$ und $B$ haben die Einheit \tonnen; $A$ ist dimensionslos. Intensitätspfad, Produktionsdaten und Emissionsprojektion müssen dieselben Scopes, Produkte und organisatorischen beziehungsweise Anlagengrenzen abdecken. Quelle, Szenarioname, Revision und Transformationsregeln des Sektorpfads sind zu dokumentieren; ein verfügbarer TPI- oder anderer wissenschaftlich begründeter Sektorpfad ist nicht automatisch für jedes Segment geeignet. Bei unterjährigen \emph{Emissionsraten} wäre zusätzlich über die Zeit zu integrieren.

$A>1$ bedeutet Überschreitung des \emph{gewählten} Referenzbudgets. Es ist keine beobachtete Unternehmenstemperatur. Produktionsprognosen, Restemissionen, Technologieannahmen und die Zuteilung globaler Budgets sind unsicher oder normativ. Ein produktionsproportionales Budget kann Wachstum belohnen; deshalb sind zusätzlich absolute sektorale Grenzen und kumulative Emissionen zu prüfen. Zielkurve, Fortschreibung bisheriger Entwicklung und technisch belegter Umsetzungspfad bleiben getrennte Szenarien. Kein Zielpfad wird als Prognose ausgegeben.

Für tatsächlich \emph{fällige absolute Emissionsziele} wird eine Zielverfehlung beispielsweise relativ zu den positiven Basisemissionen $E_{i0}$ ausgedrückt:
\begin{equation}
g_{iT}=\frac{E_{iT}^{\mathrm{beobachtet}}-E_{iT}^{\mathrm{Ziel}}}{E_{i0}}.
\end{equation}
Ein zugesagtes absolutes Reduktionsziel $r^\star$ wird dazu bei gleicher Grenze als $E_{iT}^{\mathrm{Ziel}}=(1-r^\star)E_{i0}$ geschrieben. Für \emph{Intensitätsziele} gilt separat
\begin{equation}
h_{iT}=\frac{I_{iT}^{\mathrm{beobachtet}}-I_{iT}^{\mathrm{Ziel}}}{I_{i0}}\quad\text{für }I_{i0}>0,
\end{equation}
wobei beobachtete und angestrebte Intensität dieselbe physische oder ökonomische Bezugsgröße verwenden müssen. Tonnen und Intensitäten werden nicht im selben Zähler vermischt. Damit wird auch bei einem Zielwert null nicht durch null dividiert. Bei noch nicht fälligen Zielen ist eine Zwischenpfadabweichung nur dann eine Zielverfehlung, wenn ein entsprechendes Zwischenziel tatsächlich zugesagt wurde. Eine selbst angenommene lineare Kurve ist lediglich ein Analysebenchmark.

\subsection{Capex und Anlagenbindung: der Test der Umsetzung}
Ausgeführte und bewilligte Investitionen sind näher an Umsetzung als ein fernes Versprechen, aber noch kein Wirkungsergebnis. Für positive gesamte Investitionsausgaben gilt
\begin{equation}
\kappa_i=\frac{\mathrm{Capex}_i^{\mathrm{nachweislich\ kompatibel}}}{\mathrm{Capex}_i^{\mathrm{gesamt}}}.
\end{equation}
\glqq Kompatibel\grqq{} erfordert vorab definierte technische Kriterien. Ein Marketinglabel \glqq grün\grqq{} oder bloße Taxonomie-Fähigkeit genügt nicht. Auch Taxonomie-Konformität und Kompatibilität mit jedem Net-Zero-Pfad sind nicht identisch. Ausgaben für Wartung, Expansion, Forschung und Übergangstechnologien werden getrennt klassifiziert.

Sind von $K>0$ Gesamtausgaben $L$ sicher kompatibel und $U$ sicher inkompatibel, während der Rest unklassifiziert ist, gilt unter disjunkter und vollständiger Aufteilung:
\begin{equation}
\kappa_i\in\left[\frac{L}{K},\,1-\frac{U}{K}\right].
\end{equation}
Das ist ein identifizierter Wertebereich, kein statistisches Konfidenzintervall. Unbekannte Ausgaben werden nicht pauschal als fossil oder grün etikettiert.

Die Bindung an langlebige Anlagen lässt sich szenarioabhängig als
\begin{equation}
L_i^{(s)}=\sum_{a\in\mathcal A_i}\sum_{\tau=t_0}^{T_a}Q_{a\tau}^{(s)}\,EF_{a\tau}^{(s)}
\end{equation}
modellieren: jährlicher Output mal passender Emissionsfaktor ergibt jährliche \tonnen, die zeitlich summiert werden. Laufzeit, Auslastung, Stilllegung und Umrüstung sind explizite Annahmen. $L$ ist \emph{potenzielle} zukünftige Emission unter diesen Annahmen, keine unvermeidbare Vorhersage.

\subsection{Aussagekräftige Größen jenseits des Klimas}
\small
\begin{longtable}{@{}P{0.20\linewidth}P{0.39\linewidth}P{0.35\linewidth}@{}}
\toprule
\textbf{Dimension} & \textbf{Ergebnisnahe Größe} & \textbf{Wichtige Grenze} \\
\midrule
\endfirsthead
\toprule
\textbf{Dimension} & \textbf{Ergebnisnahe Größe} & \textbf{Wichtige Grenze} \\
\midrule
\endhead
Wasser & Verbrauch nach Einzugsgebiet und Zeitraum; ergänzend $\sum_{l,m}V^{\mathrm{Verbrauch}}_{lm}CF^{\mathrm{AWARE}}_{lm}$ in knappheitsgewichteten m$^3$-Äquivalenten \cite{boulay2018}. & AWARE charakterisiert \emph{Verbrauch}, nicht pauschal Entnahme. Mengen allein ignorieren Knappheit; auch AWARE ist ein modellierter potenzieller Effekt, keine vollständige Ökosystemschadensmessung. \\
Lokale Schadstoffe & Stoff- und anlagenspezifische Freisetzungen; ergänzend EPA-RSEI unter einheitlicher Modellversion \cite{rsei}. & RSEI ist ein dimensionsloser, relativer Screeningindikator aus Toxizität, modellierter Dosis und Bevölkerung, keine Zahl verursachter Erkrankungen. US-Abdeckung ist keine globale Unternehmensbilanz. \\
Arbeitssicherheit & Todesfälle, schwere Verletzungen und DART-/Unfallraten getrennt; $200\,000N/H$ für Fälle $N$ und geleistete Stunden $H>0$. Die Basis entspricht 100 Vollzeitäquivalenten mit je 2\,000 Stunden pro Jahr \cite{bls}. & Mitarbeiter und Auftragnehmer mit passenden Zählern und Nennern erfassen; Untererfassung beachten. Abweichende Bezugsbasen anderer Raten ausdrücklich kennzeichnen. Seltene schwere Ereignisse nicht in einem Mittelwert verstecken. \\
Menschenrechte & Nachgewiesene schwere Vorfälle, betroffene Personen, Abhilfe und Wiederholungsfälle; Due-Diligence-Abdeckung separat. & Mehr Beschwerden können bessere Meldekanäle bedeuten. Null bekannte Fälle bedeutet nicht null tatsächliche Schäden. Quellenstatus und Schwere sind wichtiger als bloße Nachrichtenanzahl. \\
Biodiversität & Flächenumwandlung, Zustand und Lage betroffener Habitate; Eingriffe in sensible Gebiete. & Ein Hektar ist nicht überall ökologisch gleichwertig. Globale Standort- und Lieferkettendaten sind ein erhebliches Hindernis; keine universelle hektarbasierte Rangliste. \\
Governance & Dokumentierte Zuständigkeit, durchgesetzte Kontrollen, Abhilfe, Vergütungsmechanik und Interessenkonflikte. & Prozess- und Anreizvariablen. Ihre Wirkung auf Ergebnisse wird geprüft; ein guter Prozessscore kompensiert keinen Umwelt- oder Menschenrechtsschaden. \\
\bottomrule
\end{longtable}
\normalsize

\section{Greenwashing: von Verdacht zu überprüfbarer Aussage}
\label{sec:greenwashing}
\subsection{Nicht Nachhaltigkeitsniveau, sondern Aussage und Evidenz vergleichen}
Delmas und Burbano diskutieren organisatorische, individuelle und externe Treiber von Greenwashing; Marquis et al.\ untersuchen selektive Offenlegung \cite{delmas2011,marquis2016}. Daraus folgt nicht, dass jede unvollständige Berichterstattung absichtliche Täuschung ist. Ein legitimer enger Berichtsumfang kann transparent sein; dieselbe enge Bilanz kann durch einen umfassenden Werbeanspruch irreführend präsentiert werden.

\textbf{Prüfeinheit ist die einzelne Behauptung, nicht der moralische Gesamtcharakter einer Firma.} Eine Behauptung wird als Tupel erfasst:
\begin{equation}
C=(\text{Akteur},\text{Aussage},\text{Umfang},t_0,T,\text{Methode},\text{brutto/netto},\text{Quelle}).
\end{equation}
Dazu gehören das Originalzitat mit Fundstelle, Veröffentlichungsdatum und die Unterscheidung zwischen vergangenem Ergebnis, aktuellem Zustand, Ziel und Prognose. Ein Spartenstatement darf nicht zum Konzernversprechen umgedeutet werden. Ein Ziel kann zum Zeitpunkt seiner Ankündigung ambitioniert, aber nicht bereits widerlegt sein.

\subsection{Eine Evidenzleiter mit bewusst hohen Hürden}
\small
\begin{longtable}{@{}P{0.22\linewidth}P{0.37\linewidth}P{0.35\linewidth}@{}}
\toprule
\textbf{Status} & \textbf{Erforderlicher Befund} & \textbf{Zulässige Ausgabe} \\
\midrule
\endfirsthead
\toprule
\textbf{Status} & \textbf{Erforderlicher Befund} & \textbf{Zulässige Ausgabe} \\
\midrule
\endhead
Unzureichende Evidenz & Wesentliche Daten, Grenzen oder Originalaussage fehlen. & \glqq Nicht entscheidbar; folgende Information fehlt.\grqq{} Weder Entlastung noch Täuschungsnachweis. \\
Offenlegungslücke / Prüfhinweis & Eine relevante Kategorie oder Umsetzungsinformation ist nicht belegt. & \glqq Nicht ausreichend belegt\grqq{} oder \glqq eingeschränkte Abdeckung\grqq, nicht \glqq Unternehmen täuscht\grqq. \\
Dokumentierte Spannung & Ein konkreter Plan, eine Investition oder politische Position steht unter benannten Annahmen in Spannung zur Aussage. & \glqq Potenzielle Inkonsistenz\grqq{} mit Begründung und alternativen Erklärungen. Kein automatischer Widerspruch bei jeder fossilen Investition. \\
Quantitativer Widerspruch & Für eine vergangene oder gegenwärtige Aussage liegen vergleichbare Daten vor, deren plausible Wertebereiche die Behauptung ausschließen. & \glqq Diese abgegrenzte Aussage wird durch diese Daten widerlegt\grqq, unter Angabe von Umfang und Unsicherheit. \\
Behördlich / gerichtlich festgestellt & Eine zuständige Stelle hat einen konkreten Sachverhalt festgestellt; Datum und Verfahrensstatus sind bekannt. & Genau diesen Befund wiedergeben. Nicht von einer Entscheidung auf sämtliche Aktivitäten der Firma schließen. \\
\bottomrule
\end{longtable}
\normalsize
Eine höhere Stufe wird nicht allein durch Addition vieler schwacher Hinweise erreicht. Zehn voneinander abhängige Medienberichte über dieselbe Ursprungsquelle sind kein zehnfach unabhängiger Nachweis. Vorsatz wird durch dieses System nicht automatisch bestimmt.

\subsection{Messbare Kanäle und ihre Gegenprüfungen}
\begin{enumerate}
\item \textbf{Selektive Bilanzgrenze.} Wird aus einer Scope-1/2-Verbesserung eine konzernweite Aussage einschließlich wesentlicher Produkt- oder Lieferkettenemissionen? Gegenprüfung: Wortlaut, Scopes, relevante Kategorien und Zielabdeckung.
\item \textbf{Bilanzierung statt operativer Reduktion.} Wie viel einer gemeldeten Änderung ist mit LB/MB-Wahl, Zertifikaten oder Methodenrevisionen verbunden? Gegenprüfung: getrennte Bilanzreihen, Standardkonformität und zusätzliche Wirkung des Instruments.
\item \textbf{Verkauf oder Outsourcing.} Geht eine Bilanzverbesserung mit Anlagenübertragungen oder verlagerten Produktionsschritten einher? Gegenprüfung: Unternehmensüberleitung, konstante Anlagenkohorte und relevante Scope-3-Kategorien.
\item \textbf{Versprechen ohne Umsetzung.} Gibt es quantifizierte Zwischenziele, nachvollziehbare technische Maßnahmen und belegte Ausgaben? Gegenprüfung: Fälligkeit, genehmigte versus ausgeführte Investitionen und nachträglich geänderte Ziele.
\item \textbf{Offsets statt Bruttominderung.} Werden reale Emissionsmengen mit Kompensationskäufen verwechselt? Gegenprüfung: Bruttoinventar, registrierte und stillgelegte Credits, Projekt, Vintage, Additionalität, Baseline, Dauerhaftigkeit und Verlagerungseffekte.
\item \textbf{Widersprüchliche Einflussnahme.} Unterstützt das Unternehmen nachweislich politische Positionen, die einer eigenen spezifischen Zusage entgegenstehen? Gegenprüfung: Originalposition, Datum, Geltungsbereich und Verantwortung. Verbandsmitgliedschaft allein ist nicht gleich Zustimmung zu jeder Verbandsposition.
\end{enumerate}

Brown et al.\ finden bei 3\,574 Unternehmen mit Klimazusagen in ihrer Stichprobe bei 96\,\% mindestens einen von mehreren Risikoindikatoren, darunter häufig unvollständige Scope-3-Abdeckung \cite{brown2026}. \textbf{Das ist keine validierte Betrugsprävalenz von 96\,\%.} Ein einzelnes, sehr häufiges Merkmal kann geringe Trennschärfe besitzen. Um seine diagnostische Güte zu kennen, wären unabhängige Referenzurteile und Fehlerraten erforderlich.

Badgley et al.\ identifizieren systematische Überkreditierung in einem untersuchten kalifornischen Wald-Offset-Programm \cite{badgley2022}. Dieser Befund begründet eine Prüfung konkreter Credits; er rechtfertigt weder die Behauptung, alle Offsets seien wertlos, noch die automatische Widerlegung jedes Unternehmensclaims mit Offset-Bezug. Bilanzierungskonformität und tatsächliche Additionalität sind getrennte Fragen.

\subsection{Quantitativer Widerspruch statt willkürlichem Greenwashing-Score}
Eine Firma behaupte für einen vergleichbaren Umfang eine bereits erreichte Reduktion von mindestens $r^\star$. Sei $[\underline r,\overline r]$ der aus dokumentierten Eingabegrenzen abgeleitete Bereich der realisierten Reduktion. Dann gilt als konservative Prüfregel:
\begin{equation}
\overline r<r^\star\ \Longrightarrow\ \text{quantitativer Widerspruch unter den angegebenen Grenzen}.
\end{equation}
Liegt $r^\star$ innerhalb des Bereichs, ist der Claim nicht widerlegt. Das macht ihn nicht automatisch verifiziert. Die Grenzen dürfen nicht aus einer ungeprüften LLM-\glqq Confidence\grqq{} stammen. Für kategoriale Aussagen ist stattdessen eine explizite logische Prüfung erforderlich.

Ein \glqq Say--Do Gap\grqq{} aus der Differenz zweier beliebig normalisierter Scores ist dafür zu schwach: Aus hohem Textsentiment minus niedriger Emissionsleistung folgt kein spezifischer Widerspruch. Textanalyse kann Aussagen finden und strukturieren; sie ersetzt nicht deren sachliche Prüfung.

\section{Synthetische Gegenbeispiele: Tests gegen irreführende Kennzahlen}
\label{sec:examples}
\textbf{Alle Zahlen dieses Abschnitts sind erfundene, überprüfbare Rechenbeispiele, keine Unternehmensdaten.}
\begin{description}
\item[Preisillusion.] Emissionen bleiben bei 100\,000~\tonnen/Jahr; Umsatz steigt von 100 auf 125~Mio.~USD. Die Intensität fällt von 1\,000 auf 800~\tonnen/Mio.~USD, also um 20\,\%. Absolute Reduktion: null. Der eng formulierte Intensitätsclaim stimmt; \glqq unsere Emissionen sind um 20\,\% gesunken\grqq{} wäre damit nicht gedeckt.
\item[Anlagenverkauf.] Ein Konzern emittiert 100\,000~\tonnen/Jahr und verkauft eine Anlage mit 40\,000~\tonnen/Jahr. Beide verbleibenden und verkauften Anlagen produzieren unverändert. Konzerninventar danach: 60\,000; konstante Anlagenkohorte: weiterhin 100\,000. Die Bilanz sinkt um 40\,\%, ohne physische Reduktion dieser Kohorte.
\item[Strombilanz.] Scope~2 LB beträgt in beiden Jahren 100\,000~\tonnen. MB sinkt durch Beschaffungsinstrumente von 100\,000 auf 10\,000. Das belegt 90\,\% MB-Bilanzreduktion, nicht für sich genommen 90\,\% zusätzliche Klimawirkung. Deren Höhe bleibt ohne Interventionsanalyse offen.
\item[Scope-Verlagerung.] Scope~1/2 sinkt von 100 auf 70 Einheiten, während vergleichbar erfasstes Scope~3 von 900 auf 1\,100 steigt. Die partielle Bilanz fällt um 30\,\%; die hier vollständig angenommene Gesamtbilanz steigt von 1\,000 auf 1\,170, also um 17\,\%.
\item[Effizienz mit Wachstum.] Die physische Intensität sinkt um 10\,\%, der Output steigt um 30\,\%. Absolute Emissionen steigen um $0{,}9\times1{,}3-1=17\,\%$. Effizienzfortschritt und höhere Gesamtbelastung sind gleichzeitig wahr.
\item[Gleicher Endpunkt, anderes Budget.] Zwei Pfade enden im vierten Jahr bei null. Pfad A emittiert jährlich $(80,60,40,0)$, Pfad B $(100,100,100,0)$ Tausend~\tonnen. Kumulativ ergeben sich 180 gegenüber 300 Tausend~\tonnen. Ein gemeinsames Net-Zero-Jahr verschleiert 120 Tausend~\tonnen{} Unterschied.
\item[Datenlücke.] Unternehmen A berichtet 100 Einheiten Scope~1/2 und 900 Scope~3. B berichtet 120 Scope~1/2 und kein Scope~3. Ein Gesamtvergleich \glqq B ist besser\grqq{} ist unzulässig. Nur die vergleichbare Teilbilanz kann verglichen werden; die Gesamtordnung bleibt offen.
\end{description}
Diese Beispiele sind wertvoller als eine bloße Demo mit eindrucksvollen Ranglisten: Ein Bewertungsverfahren sollte nachweisbar verhindern, dass solche Effekte als unbedingter Nachhaltigkeitsfortschritt ausgegeben werden.

\section{Eigene Methodik: Claim--Evidence--Outcome Audit}
\label{sec:design}
\subsection{Architektur und Datenmodell}
Der vorgeschlagene Beitrag besteht aus drei verknüpften Ebenen:
\begin{enumerate}
\item \textbf{Outcome Ledger:} versionierte Kennzahlen pro Unternehmen, Anlage oder Segment; Einheit, Zeitraum, Scopes und Vergleichbarkeit sind Pflichtfelder.
\item \textbf{Claim Ledger:} konkrete Originalaussagen mit Quelle, Datum, Umfang und Typ. Ergebnisbehauptung, Ziel und Prognose werden getrennt.
\item \textbf{Evidence Graph:} Beziehungen \glqq stützt\grqq, \glqq widerspricht\grqq{} und \glqq reicht nicht zur Prüfung\grqq{} zwischen Aussagen, Messungen und Dokumenten. Quellenabhängigkeiten und alternative Erklärungen bleiben sichtbar.
\end{enumerate}
Das ist keine Behauptung eines neuen Grundlagenalgorithmus. Der wissenschaftliche Beitrag wäre die nachprüfbare Operationalisierung samt Benchmark, Grenzbehandlung und empirisch gemessener Verbesserung gegenüber einfacheren Verfahren.

\textbf{KI-Rolle.} Ein Sprachmodell kann Fundstellen und Tabellen extrahieren und strukturierte Kandidaten liefern. Es darf fehlende Zahlen nicht ergänzen oder Gesamturteile frei generieren. Einheitenprüfung, Berechnung und Entscheidungsregeln erfolgen deterministisch. Eine wörtliche Fundstelle belegt lediglich, dass die Quelle etwas sagt; sie beweist noch nicht dessen Wahrheit. Übernommene Modellwerte werden als solche markiert und nicht zu beobachteten Fakten umetikettiert.

\subsection{Datenqualität nicht mit Unternehmensleistung verrechnen}
Jeder Datensatz erhält getrennte Merkmale für Abdeckung, zeitliche Aktualität, unabhängige Prüfung, Unsicherheit, Quellenabhängigkeit und Grenzkonsistenz. Kein zusammenfassender \glqq Datenqualitätspunkt\grqq{} wird zur Belohnung des Unternehmens auf die Umweltleistung addiert. Gute Berichterstattung kann große Schäden sichtbar machen und verdient nicht deshalb ein schlechteres Leistungsurteil als fehlende Berichterstattung.

Fehlende Werte können nicht zufällig fehlen: Branche, Komplexität und Anreize beeinflussen Offenlegung. Deshalb ist ein Complete-Case-Ranking ebenfalls potenziell selektiv. Das System berichtet die Abdeckung nach Sektor und Unternehmensgröße; Imputationen sind allenfalls ein separat gekennzeichnetes Sensitivitätsszenario, nicht der Standardwert.

\subsection{Robuste Vergleiche und partielle Ordnung}
Auf kompatiblen, nachteilsorientierten Kriterien $d_{ik}$ mit belastbaren Eingabeintervallen dominiert A die Firma B konservativ, wenn
\begin{equation}
\overline d_{Ak}\leq\underline d_{Bk}\quad\text{für alle berücksichtigten }k
\end{equation}
und die Ungleichung für mindestens ein Kriterium strikt ist. Andernfalls kann der Vergleich offenbleiben. Diese strenge Regel kann viele unentscheidbare Fälle erzeugen; das ist informativer als erzwungene Scheingenauigkeit. Bei abhängigen Eingaben sind gemeinsame zulässige Szenarien aussagekräftiger als unabhängig kombinierte Randintervalle.

Wo ein Nutzer eine Rangliste benötigt, werden Richtung, Skalierung und zulässige Gewichte vorab festgelegt. Mit normalisierten Kostenmerkmalen $z_{ik}^{(s)}$ und Szenarien $s$ gilt beispielsweise
\begin{equation}
S_i^{(s)}=\sum_k w_k^{(s)}z_{ik}^{(s)},\qquad w_k^{(s)}\geq0,\quad\sum_k w_k^{(s)}=1.
\end{equation}
Die Szenarien variieren \emph{nicht nur Gewichte}, sondern auch plausible Grenz- und Messannahmen. Die Methodik orientiert sich an der Idee stochastischer multikriterieller Akzeptabilitätsanalyse \cite{smaa1998}. Ausgegeben werden Rangbereiche und der Anteil spezifizierter Szenarien, unter denen ein Vergleich gilt. Diese Anteile sind ohne kalibrierte Verteilungsannahmen keine objektiven Nachhaltigkeitswahrscheinlichkeiten. Ein robustes Ranking mit unvollständigen falschen Daten bleibt falsch.

\subsection{Keine vollständige Kompensation schwerer Schäden}
Vor einer optionalen Aggregation werden dokumentierte schwere Vorfälle und nicht erfüllte wesentliche Mindestanforderungen separat angezeigt. Ob sie zum Ausschluss führen, ist eine ausdrücklich gewählte Nutzungsregel. Unbekannte Informationen gelten dabei nicht automatisch als Regelverletzung. Ökologische, soziale und finanzielle Zielkonflikte bleiben sichtbar; das System verkauft einen normativen Gewichtungsentscheid nicht als Naturgesetz.

\section{Wie die Aussagekraft empirisch überprüft werden müsste}
\label{sec:validation}
\subsection{Forschungsfrage 1: Bringt ESG zusätzliche Prognoseinformation?}
Die relevante Prüfung lautet nicht, ob ein Gesamtscore mit seinen eigenen Komponenten korreliert. Das wäre teilweise konstruktionsbedingt. Zu testen ist, ob der Score \emph{unabhängige spätere Ergebnisse} besser vorhersagt als einfache Baselines.

Verglichen werden beispielsweise:
\begin{itemize}
\item Baseline B0: Branche, Größe, Output und vergangene Ergebniswerte;
\item B1: B0 plus zeitgenössisch verfügbarer ESG- beziehungsweise E-Score;
\item B2: B0 plus konkrete Maßnahmen, Investitionen und Abdeckungsinformationen;
\item B3: B2 plus vorab erfasste, belegte Claim-Inkonsistenzen.
\end{itemize}
Zielvariablen sind vergleichbare zukünftige physische Intensität, Emissionen auf konstanter Anlagenbasis oder unabhängig dokumentierte schwere Ereignisse. Absolute Emissionen allein hängen stark von Unternehmensgröße und Geschäftsmodell ab. Verstöße hängen zudem von Kontrolldichte, Rechtssystem und Entdeckungswahrscheinlichkeit ab; sie sind keine perfekte Schadens-Ground-Truth.

Ein schematisches Prognosemodell ist
\begin{equation}
Y_{i,t+h}=\alpha_{\mathrm{Sektor}(i)}+\gamma_t+\rho Y_{it}+\beta R_{it}+\delta^\top X_{it}+\varepsilon_{i,t+h}.
\end{equation}
Die Spezifikation ist eine Ausgangshypothese, kein bereits geschätztes Modell. Zeitliche Validierung, Modellkomplexität und ausreichend unabhängige Beobachtungen sind wichtiger als ein ausgewählter signifikanter Koeffizient. Bei Panelinferenz sind Abhängigkeiten innerhalb von Unternehmen zu berücksichtigen; dynamische Fixed-Effects-Modelle mit kurzen Zeitreihen haben eigene Verzerrungsprobleme. Ein prädiktiver Zugewinn beweist nicht, dass das Rating eine physische Veränderung verursacht.

\subsection{Forschungsfrage 2: Verbessert Aussageprüfung die Erkennung von Widersprüchen?}
Es wird ein Referenzset aus konkreten Claims erstellt, stratifiziert nach Sektor und Aussageart. Zwei voneinander unabhängige Prüfer ordnen Fälle anhand derselben Anleitung ein; Konflikte werden dokumentiert und adjudiziert. Prozentuale Übereinstimmung und ein geeignetes zufallskorrigiertes Maß werden berichtet. Ein hoher Kappa-Wert ist weder garantiert noch ein Wahrheitsbeweis.

Getrennte Messgrößen sind:
\begin{description}
\item[Extraktionsgüte:] Übereinstimmung von Wert, Einheit, Jahr, Scope, Grenze und Originalfundstelle. Eine richtige Zahl mit falschem Scope ist ein Fehler.
\item[Widerspruchspräzision:] Anteil bestätigter Widersprüche unter den als widersprüchlich ausgegebenen Fällen, mit Unsicherheitsangabe.
\item[Recall:] Anteil gefundener Widersprüche im unabhängig zusammengestellten Referenzset, nicht nur in bereits vom Modell ausgewählten Fällen.
\item[Abdeckung:] Anteil der Claims, zu denen das System ein hinreichend belegtes Urteil abgibt; Präzision und Enthaltungsquote werden gemeinsam ausgewiesen.
\item[Fehlerstruktur:] Fehler nach Branche, Claim-Typ und Kanal, nicht nur eine Gesamt-Accuracy bei möglicherweise stark ungleichen Klassen.
\end{description}
Die Referenzlabels lauten \glqq dokumentierter Widerspruch\grqq, \glqq ausreichende Evidenz\grqq{} oder \glqq unzureichende Evidenz\grqq, nicht automatisch \glqq Betrug\grqq. Synthetische Fixtures aus Abschnitt~\ref{sec:examples} prüfen Rechenlogik; sie ersetzen keine Validierung an echten, bisher ungesehenen Fällen.

\subsection{Forschungsfrage 3: Unterscheidet das Modell echte von bilanzieller Verbesserung?}
Für Anlagenverkäufe werden Unternehmensinventar, Rating und dieselben Anlagen vor und nach dem Verkauf verfolgt. Eine Ereignisstudie mit geeigneten Vergleichsanlagen kann Mechanismen untersuchen. Selektive Verkäufe, unterschiedliche Käufer, Produktionseinbrüche und regulatorische Änderungen verhindern jedoch eine naive Kausalinterpretation. Eine Difference-in-Differences-Analyse benötigt nachvollziehbare Identifikationsannahmen; die bloße Bezeichnung \glqq DiD\grqq{} schafft diese nicht.

Zusätzlich wird mittels Ablationen geprüft, welchen Nutzen Grenzüberleitung, LB/MB-Trennung und Claim-Strukturierung jeweils bringen. Eine interessante technische Kennzahl ist die Rate \emph{falsch behaupteter physischer Verbesserungen} gegenüber einer einfachen Umsatzintensitäts- oder ESG-Rangliste. Die Tests müssen auch echte Verbesserungen enthalten, damit ein System nicht durch permanente Skepsis scheinbar erfolgreich wird.

\subsection{Schutz vor zeitlicher und inhaltlicher Datenleckage}
\begin{enumerate}
\item \textbf{Beobachtungsjahr, Veröffentlichungsdatum und Abrufversion getrennt speichern.} Ein Bericht für 2024, veröffentlicht 2025, war 2024 nicht vollständig bekannt.
\item \textbf{Zeitlich getrennte Entwicklung und Prüfung.} Schwellen und Prompts nicht nach Sichtung der Testfälle anpassen; zusammengehörige Dokumente eines Falls nicht über Trainings- und Testmengen verteilen.
\item \textbf{Historische Indexmitgliedschaft beachten.} Heutige S\&P-500-Mitglieder als historischen Investitionsbestand zu verwenden erzeugt Survivorship-Bias. Ein heutiger Querschnitt kann dagegen ausdrücklich das aktuelle Universum untersuchen.
\item \textbf{Keine zukünftigen Ziele rückwirkend benutzen.} Gelöschte, revidierte und verfehlte Ziele bleiben im historischen Datensatz. Noch nicht fällige Ziele werden nicht als Erfolge oder Misserfolge etikettiert.
\item \textbf{Keine zirkuläre Validierung.} Ein aus SBTi abgeleiteter Glaubwürdigkeitsscore darf nicht mit SBTi-Status als unabhängigem Wahrheitslabel geprüft werden. Dasselbe gilt für zwei Anbieter, die denselben Unternehmensbericht übernehmen.
\item \textbf{Vergleiche auf gleicher Stichprobe.} Ergebnisse für ein kleines gut dokumentiertes Teiluniversum dürfen nicht als flächendeckender Nachweis für alle Unternehmen ausgegeben werden.
\end{enumerate}
Ein Hackathon kann diese Infrastruktur und eine kleine manuell geprüfte Stichprobe demonstrieren. Er kann ohne passendes Datenpanel keine belastbare globale Kausalstudie ersetzen. Es werden vor Durchführung weder Modellgenauigkeit noch Signifikanz oder Überrendite behauptet.

\section{Datenquellen, Abdeckung und realistische Umsetzung}
\label{sec:data}
\subsection{Was in dieser Recherche tatsächlich geprüft wurde}
Die S\&P-500-Referenz-CSV wurde erfolgreich abgerufen; sie enthält 503 Wertpapierzeilen. Mehrere Aktienklassen sind nicht mehrere unabhängige Unternehmen \cite{constituents}. CIK, LEI und gegebenenfalls ISIN werden für Unternehmensidentitäten verwendet; Namensähnlichkeit allein genügt nicht. Der Abgleich zwischen sämtlichen Quellen wurde noch nicht ausgeführt, daher wird keine erreichte Unternehmensabdeckung behauptet.

Die Excel-Dateien des SBTi-Dashboards und des Climate-Action-100+-Benchmarks wurden erfolgreich heruntergeladen und auf Tabellenstruktur geprüft \cite{sbti,ca100}. Das SBTi-Exportformat enthält unter anderem Unternehmensidentifikatoren, Zielstatus, Zieljahre und Zieltexte; es ist keine vollständige Emissionshistorie. Die geprüfte CA100+-Ausgabe betrifft die 2025er Bewertungen von 164 globalen Fokusunternehmen, nicht 164 nachgewiesene S\&P-500-Mitglieder. Die einzelnen Bewertungsteile verwenden unterschiedliche Datenstichtage.

\small
\begin{longtable}{@{}P{0.22\linewidth}P{0.36\linewidth}P{0.36\linewidth}@{}}
\toprule
\textbf{Quelle} & \textbf{Rolle} & \textbf{Grenze / Zugangsstatus} \\
\midrule
\endfirsthead
\toprule
\textbf{Quelle} & \textbf{Rolle} & \textbf{Grenze / Zugangsstatus} \\
\midrule
\endhead
Unternehmensberichte & Inventare, Produktionsdaten, Originalclaims, Investitionsbelege. & Selbstberichte; Validierung, Einheitentransformation und Rechte an Dokumenten erforderlich. Verifizierung eines Berichts ist nicht zwingend Verifizierung aller Claims. \\
SBTi \cite{sbti} & Zielstatus und abgegrenzte validierte Ziele. & \glqq Committed\grqq{} ist nicht \glqq Targets set\grqq. Keine Zertifizierung der gesamten Firma oder ihrer tatsächlichen künftigen Zielerreichung. \\
CA100+ \cite{ca100} & Strukturierte Beurteilungen von Übergangsplanung und Umsetzung. & Erfolgreich geprüfter Excel-Export, begrenztes Universum; Nutzungsbedingungen einzelner Datenpartner unterscheiden sich. \\
Net Zero Tracker \cite{nzt2025} & Zielabdeckung, Zwischenziele, Plan- und Offsetmerkmale. & 2025er Zenodo-Snapshot als CSV mit CC BY 4.0 dokumentiert. Direkter Dateiabruf in dieser Umgebung: HTTP 403. Manuell zugänglicher Download ist vor Einsatz zu bestätigen. \\
TPI \cite{tpi,tpiterms} & Sektorpfade und Methodik; möglicher externer Vergleich. & Nutzungsbedingungen enthalten Genehmigungsanforderungen u.\,a. für automatisierte Nutzung und Fonds-/Indexkonstruktion sowie Grenzen der Weitergabe. Öffentlich sichtbar ist nicht uneingeschränkt frei nutzbar. \\
WBA \cite{wba} & Breitere Klima-, Sozial- und Unternehmensindikatoren. & Datenportal benötigt Registrierung. API-Freigabe laut Dokumentation möglicherweise erst nach mehreren Tagen. Kein kritischer Pfad für den Hackathon. \\
EPA GHGRP / RSEI \cite{epaghg,rsei} & Anlagenbezogene US-Emissionen und lokale schadstoffbezogene Screeningdaten. & Abdeckungs-, Schwellen- und Modellgrenzen; keine vollständige globale Konzernbilanz. Anlagen-/Eigentümerabgleich erforderlich. \\
Climate TRACE \cite{climatetrace} & Ergänzende anlagenbezogene Schätzungen und Ownership-Informationen. & Abdeckungs- und Modellunterschiede; nicht als vollständige Scope-1/2/3-Inventur behandeln und nicht doppelt mit denselben EPA-Anlagen addieren. \\
SEC CompanyFacts \cite{sec} & Finanzielle Nenner und Unternehmensidentifikation. & XBRL-Tags, Jahres-/Quartalswerte und Restatements prüfen. Keine vorausgesetzte vollständige GHG-Datenbank; EBITDA nicht als universell vorhandenen standardisierten Tag annehmen. \\
\bottomrule
\end{longtable}
\normalsize
Die Zugangsprüfung ersetzt keine Lizenzprüfung für Weiterveröffentlichung oder kommerzielle Nutzung. Ein vollständiger, frei zugänglicher und geprüfter Scope-1/2/3-Zeitreihendatensatz für sämtliche S\&P-500-Unternehmen wurde in dieser Recherche nicht gefunden; daraus folgt kein Beweis, dass ein solcher Datensatz nirgends existiert.

\subsection{Stufenweises Vorgehen statt erfundener Vollständigkeit}
\begin{enumerate}
\item \textbf{Universum und Coverage Audit.} Unternehmen und Aktienklassen trennen; je Kennzahl und Branche erreichbare Abdeckung, Zeitstand und Rechte bestimmen.
\item \textbf{Kleiner, tiefer Benchmark.} Beispielsweise 20--30 geeignete Unternehmen als Umfangsziel, darunter mehrere emissionsintensive und weniger emissionsintensive Branchen. Auswahlregeln vorher dokumentieren; keine erreichte Stichprobengröße behaupten.
\item \textbf{Verifizierte Fallpaare.} Zwei bis drei aussagekräftige Vergleiche mit Originalclaims, Datenüberleitung und Gegenbelegen; mindestens ein echter Fortschritt und ein unentscheidbarer Fall.
\item \textbf{Deterministische Tests.} Die Gegenbeispiele aus Abschnitt~\ref{sec:examples}, Nenner-null-Fälle, Einheitentausch, Aktienklassenduplikate und Scope-Verwechslungen als Testsuite.
\item \textbf{Skalierung nur bei Datenreife.} Weitere Unternehmen hinzufügen, ohne Abdeckungsanforderungen zu lockern. Nicht bewertbare Unternehmen bleiben sichtbar statt aus der Oberfläche zu verschwinden.
\end{enumerate}
Sektorpfade sind für vergleichbare Technologien beziehungsweise Segmente auszuwählen, nicht nur anhand der elf groben GICS-Sektoren. Ein Konzern mit mehreren Geschäftsmodellen wird nicht durch einen einzigen scheinpräzisen Sektorpfad erklärt.

\section{Net-Zero-Fonds: Zurechnung, Ausrichtung und tatsächlicher Einfluss}
\label{sec:portfolio}
\subsection{Die finanzierte Emission ist kein kausaler Fußabdruck des Investors}
Für gelistete Aktien und Unternehmensanleihen wird die PCAF-Zurechnung schematisch als
\begin{equation}
FE_p=\sum_i\frac{V_{pi}}{\mathrm{EVIC}_i}\,E_i
\end{equation}
dargestellt. $V_{pi}$ ist der gehaltene Wert, EVIC der nach Standard definierte Unternehmenswert einschließlich Cash und $E_i$ die abgegrenzte jährliche Emissionsmenge \cite{pcaf2025}. Für diese Anlageklasse umfasst EVIC die Marktkapitalisierung der Stamm- und relevanten Vorzugsaktien, den Buchwert der Gesamtschulden sowie nicht beherrschende Anteile (minority/non-controlling interests); Cash wird nicht wie bei einer üblichen Netto-Unternehmenswertkennzahl abgezogen. Fehlende oder nichtpositive Nenner werden nicht durch willkürliche Ersatzwerte repariert. Die konkrete Standardfassung, Stichtage und Scope-Berichterstattung sind anzuwenden.

\textbf{Synthetisches Beispiel:} Bei konstantem gehaltenem Wert $V=1$~Mio.~USD und konstanten Emissionen $E=100\,000$~\tonnen{} sinkt die Zurechnung von 1\,000 auf 500~\tonnen, wenn EVIC von 100 auf 200~Mio.~USD steigt. Das ist eine Bewertungsänderung, keine operative Reduktion. Bei tatsächlich festen Stückzahlen kann sich $V$ mitbewegen; deshalb müssen Zähler, Nenner und Handelsänderungen separat übergeleitet werden.

Fraser und Fiedler zeigen, dass große Reduktionen finanzierter Emissionen mit steigenden physischen Emissionen vereinbar sind \cite{fraser2023}. PCAF sieht inzwischen ergänzende, separat berichtete vermiedene Emissionen und vorausschauende Größen vor; sie sind nicht einfach Teil der Kerninventur \cite{pcaf2025}.

\subsection{Unternehmenswirkung und Investoreneinfluss trennen}
Vermiedene Emissionen einer Maßnahme wären bei gleicher funktionaler Leistung und Systemgrenze
\begin{equation}
AE^{(s)}=E_{\mathrm{ohne\ Maßnahme}}^{(s)}-E_{\mathrm{mit\ Maßnahme}}^{(s)}.
\end{equation}
Der Gegenfakt ist unbeobachtet und muss begründet werden. Ein grünes Produkt kann andere Produkte verdrängen, zusätzliche Nachfrage erzeugen oder bereits existierende saubere Lösungen ersetzen. Deshalb ist \glqq grüne Umsätze\grqq{} kein direktes Maß zusätzlicher Klimawirkung. Vermiedene Emissionen werden nicht zur rechnerischen Löschung der eigenen Bruttoinventur verwendet.

Noch eine Ebene schwieriger ist die Frage, welchen Anteil der Investor an dieser Änderung verursacht. Der Kauf bestehender Aktien überweist dem Unternehmen nicht unmittelbar neues Projektkapital. Kölbel et al.\ finden in ihrem Review relativ stärkere Evidenz für Shareholder Engagement als für reine Kapitalallokation \cite{koelbel2020}. Hartzmark und Shue diskutieren in einem Working Paper einen gegenläufigen Kanal: Höhere Finanzierungskosten emissionsintensiver Firmen können deren Umweltleistung verschlechtern \cite{hartzmark2026}. Dies ist kein allgemeiner Beweis gegen Divestment oder für jedes braune Unternehmen, sondern ein Grund, den Wirkungsmechanismus statt einer mechanischen \glqq braun raus, grün rein\grqq{}-Regel zu prüfen.

\subsection{Ein ehrlicher Optimierungsauftrag für 1 Mrd. USD}
Ohne glaubwürdige zusätzliche Abatement-Kosten und Finanzierungseffekte kann das System nicht seriös \glqq maximal vermiedene Tonnen pro investiertem Dollar\grqq{} optimieren. Die Heuristik \glqq heutige Emissionen geteilt durch Marktkapitalisierung\grqq{} ist dafür ungeeignet: Sie verwechselt Zurechnung beziehungsweise Exposition mit beeinflussbarem Reduktionspotenzial.

Für einen Aktien-Demonstrator wird stattdessen ein begrenztes Mandat definiert: \textbf{robuste Transition-Ausrichtung bei kontrollierter Konzentration und offengelegten Datenlücken}. Als konkret berechenbare, dimensionslose Transition-Kosten wird die positive relative Überschreitung aus Gleichung~\ref{eq:budget} verwendet:
\begin{equation}
d_i^{(s)}=\max\!\left(0,\frac{C_i^{(s)}}{B_i^{(s)}}-1\right),\qquad B_i^{(s)}>0.
\end{equation}
Unternehmen ohne hinreichend begründbares Budget und Emissionsszenario gehören nicht zum bewertbaren Optimierungsuniversum; ihre Nichtberücksichtigung und der dadurch mögliche Selektionsbias werden offengelegt. Ein bloß angekündigter Zielpfad darf das operative Emissionsszenario nicht ersetzen. Der Kostenwert misst relative Pfadabweichung, nicht beeinflussbare Tonnen pro Dollar; heutige absolute Belastung und kumulative Mengen bleiben eigene Ausgaben. Die Optimierung lautet dann:
\begin{align}
\min_{w,z}\quad &z\\
\text{unter}\quad &\sum_i w_i d_i^{(s)}\leq z\quad\forall s,\\
&\sum_i w_i=1,\qquad 0\leq w_i\leq w_{\max},\\
&\left|\sum_{i\in k}w_i-b_k\right|\leq\delta_k\quad\forall k.
\end{align}
$w$ sind Gewichte, $b_k$ gewünschte Sektoranteile und $s$ veröffentlichte Mess-/Transitionsszenarien. Bei verlässlichen, zeitlich passenden Kursdaten kann zusätzlich $(w-b)^\top\Sigma(w-b)\leq\tau^2$ als Aktivrisikogrenze dienen. Ohne solche Daten wird eine Gewichtsabweichung nicht als geschätztes Marktrisiko ausgegeben. Unlösbare Nebenbedingungen werden gemeldet, nicht still gelockert. Dollarallokationen sind $10^9w_i$.

Das Ergebnis ist ein \emph{Ausrichtungsportfolio}, kein kausal optimales Klimawirkungsportfolio. Verglichen werden eine einfache Baseline auf demselben Universum, ein Low-Carbon-Ansatz und die robuste Transition-Variante. Zu berichten sind Scope-spezifische Exposition, Datenabdeckung, Konzentration, Sektorabweichung und kumulative Szenariopfade. Historische Renditen werden nicht als sichere Erträge im plötzlich geänderten politischen Regime extrapoliert; nachhaltige Anlagen können im Gleichgewicht niedrigere erwartete Renditen haben \cite{pastor2021}.

\textbf{Das Wort \glqq schnellstmöglich\grqq{} ernst nehmen:} Kurzfristige Umsetzung und frühe kumulative Reduktionen zählen, nicht nur ein Zieljahr 2050. Globales Commitment beseitigt technische und infrastrukturelle Engpässe nicht sofort. Bei freiem Anlageklassenmandat wären zusätzliche Projektfinanzierung und vertragliche Meilensteine direktere Einflusskanäle, benötigen aber ein anderes Daten- und Risikomodell.

\section{Empfehlung für ETHack und wissenschaftliche Schlussfolgerung}
\subsection{Der eigentliche Innovationskern}
Das überzeugendere Projekt lautet nicht \glqq Wir wissen, welches Unternehmen nachhaltig ist\grqq, sondern:
\begin{quote}
\textbf{Wir prüfen, welche behaupteten Nachhaltigkeitsfortschritte durch vergleichbare Daten gedeckt sind, welche nur aus Bilanzierung entstehen und wo keine belastbare Entscheidung möglich ist.}
\end{quote}
Das System kann damit die fünf Kriterien des Decks adressieren:
\begin{description}
\item[Impact:] konkrete Fehlentscheidungen und unbelegte Wirkungsclaims vermeiden; Nutzer wären beispielsweise Investmentteams, Analysten oder zivilgesellschaftliche Prüfer.
\item[Innovation:] Claim-genaue Widerspruchsprüfung mit Bilanzüberleitung, Quellenabhängigkeit und expliziter Nichtentscheidbarkeit statt generischer LLM-Zusammenfassung.
\item[Technical Execution:] funktionierende Extraktion, korrekte Einheiten, verknüpfte Belege, adversarische Tests und reproduzierbare Szenarien.
\item[Feasibility:] enger, verifizierter Datensatz; keine vorausgesetzten proprietären Vollabdeckungen oder kurzfristig nicht erhältlichen API-Freigaben.
\item[Presentation:] ein realer Vorher-nachher-Claim mit offener Belegkette, ein Gegenbeispiel und ein klarer Entscheidungseffekt. Keine unvalidierten Erfolgszahlen.
\end{description}
Für die dreiminütige Präsentation sollte ein Fall im Mittelpunkt stehen, nicht das gesamte Kennzahlensystem. Die vorgeschriebene Präsentationsdatei ist laut Deck im Format \texttt{.ppt} einzureichen; dieser Bericht ist ein ergänzendes wissenschaftliches Arbeitsdokument, kein Ersatz für die Folien \cite{deck}.

\subsection{Antwort auf die Ausgangsfrage}
\textbf{Ein ESG-Score ist weder wertlos noch ein Nachhaltigkeitsbeweis.} Seine Aussagekraft hängt von Konstrukt, Datenbasis, Anwendungszweck und externer Validierung ab. Ergebnisnahe Größen sind geeigneter zur Prüfung realer Leistung, bleiben aber von Grenzen, Standort, Aktivität und Datenqualität abhängig. Ein Unternehmen kann transparent schlechte Ergebnisse berichten oder gute Ergebnisse übertrieben darstellen. Nachhaltigkeitsleistung und Glaubwürdigkeit der Kommunikation müssen deshalb getrennt untersucht werden.

\textbf{Greenwashing ist nicht zuverlässig durch eine weitere Blackbox-Zahl erkennbar.} Der tragfähige Weg ist eine dokumentierte Kette von Originalbehauptung über passende Messung zur begrenzten Schlussfolgerung. Ein aussagekräftiger Output darf \glqq nicht entscheidbar\grqq{} lauten. Erst diese Zurückhaltung ermöglicht einen Vergleich, der wissenschaftlich mehr leistet als eine attraktive Rangliste.

\appendix
\section{Minimaler Prüfdatensatz und Freigabekriterien}
\label{app:schema}
\textbf{Pro Kennzahl:} Unternehmens-ID und gegebenenfalls Anlagen-/Segment-ID; Kennzahl; Wert; Einheit; Menge oder Rate einschließlich Bezugsdauer; Zeitraum; Veröffentlichungsdatum; Datenversion; Quelle und Seite; Scope/Kategorie; LB/MB; organisatorische Grenze; gemessen/berichtet/modelliert; geprüfter Umfang; fehlende Kategorien; Überleitung gegenüber dem Basisjahr.

\textbf{Pro Claim:} Originaltext; Originalquelle; Datum; Akteur; Gegenstand; Ergebnis/Ziel/Prognose; Basis- und Zieljahr; Scopes; Bilanzgrenze; brutto/netto; erforderliche Belege; stützende und entgegenstehende Evidenz; Quellenabhängigkeiten; alternative Erklärungen; Status; Prüferentscheidung.

\textbf{Vor Freigabe eines Vergleichs:}
\begin{enumerate}
\item Ist das Messobjekt finanzielle Resilienz, Unternehmenswirkung oder Investoreneinfluss?
\item Sind Zeitraum, Scopes, Einheiten und Grenzen kompatibel?
\item Sind Datenlücken, Revisionen und modellierte Werte sichtbar?
\item Bleibt der Befund bei plausiblen Gegenannahmen bestehen?
\item Werden nicht fällige Ziele und echte Fortschritte fair behandelt?
\item Ist die Schlussfolgerung enger oder höchstens so weit wie die Belege?
\item Sind Prüfqualität, Datenabdeckung und Portfoliowirkung nur dann quantifiziert, wenn tatsächlich gemessen?
\end{enumerate}

\newpage
\begin{thebibliography}{99}
\small
\raggedright
\bibitem{deck} ETHack (2026). \emph{ETHack 2026 -- Opening Ceremony}. Intro-Deck, insbesondere Folien 8--10 und 16. \url{https://canva.link/ethack2026-intro}. Abgerufen am 12.09.2026.

\bibitem{msci} MSCI (o.\,J.). \emph{ESG Ratings: Assess companies on their financially relevant sustainability risks and opportunities}. Offizielle Produkt- und Methodikbeschreibung. \url{https://www.msci.com/data-and-analytics/sustainability-solutions/esg-ratings}. Abgerufen am 12.09.2026.

\bibitem{sustainalytics} Morningstar Sustainalytics (o.\,J.). \emph{ESG Risk Ratings -- Methodology Abstract}. Offizielle Methodikbeschreibung. \url{https://www.sustainalytics.com/docs/default-source/default-document-library/sustainalytics_esg-risk-ratings_methodology-abstract.pdf}. Recherchestand 12.09.2026.

\bibitem{berg2022} Berg, F., Kölbel, J. F. und Rigobon, R. (2022). Aggregate Confusion: The Divergence of ESG Ratings. \emph{Review of Finance}, 26(6), 1315--1344. \doi{10.1093/rof/rfac033}.

\bibitem{christensen2022} Christensen, D. M., Serafeim, G. und Sikochi, A. (2022). Why Is Corporate Virtue in the Eye of the Beholder? The Case of ESG Ratings. \emph{The Accounting Review}, 97(1), 147--175. \doi{10.2308/TAR-2019-0506}.

\bibitem{drempetic2020} Drempetic, S., Klein, C. und Zwergel, B. (2020). The Influence of Firm Size on the ESG Score: Corporate Sustainability Ratings Under Review. \emph{Journal of Business Ethics}, 167, 333--360. Online zuerst 2019. \doi{10.1007/s10551-019-04164-1}.

\bibitem{chatterji2009} Chatterji, A. K., Levine, D. I. und Toffel, M. W. (2009). How Well Do Social Ratings Actually Measure Corporate Social Responsibility? \emph{Journal of Economics \& Management Strategy}, 18(1), 125--169. \doi{10.1111/j.1530-9134.2009.00210.x}.

\bibitem{raghunandan2022} Raghunandan, A. und Rajgopal, S. (2022). Do ESG funds make stakeholder-friendly investments? \emph{Review of Accounting Studies}, 27, 822--863. \doi{10.1007/s11142-022-09693-1}.

\bibitem{treepongkaruna2024} Treepongkaruna, S., Au Yong, H. H., Thomsen, S. und Kyaw, K. (2024). Greenwashing, Carbon Emission, and ESG. \emph{Business Strategy and the Environment}, 33(8), 8526--8539. \doi{10.1002/bse.3929}.

\bibitem{khan2016} Khan, M., Serafeim, G. und Yoon, A. (2016). Corporate Sustainability: First Evidence on Materiality. \emph{The Accounting Review}, 91(6), 1697--1724. \doi{10.2308/accr-51383}.

\bibitem{berg2021} Berg, F., Fabisik, K. und Sautner, Z. (2021). \emph{Is History Repeating Itself? The (Un)Predictable Past of ESG Ratings}. Working Paper, ECGI Finance Working Paper 708/2020; hier die 2021er Fassung. \doi{10.2139/ssrn.3722087}.

\bibitem{edmans2023} Edmans, A. (2023). The end of ESG. \emph{Financial Management}, 52(1), 3--17. Online zuerst 2022. \doi{10.1111/fima.12413}.

\bibitem{busch2022} Busch, T., Johnson, M. und Pioch, T. (2022). Corporate carbon performance data: Quo vadis? \emph{Journal of Industrial Ecology}, 26(1), 350--363. Online zuerst 2020. \doi{10.1111/jiec.13008}.

\bibitem{klaassen2021} Klaaßen, L. und Stoll, C. (2021). Harmonizing corporate carbon footprints. \emph{Nature Communications}, 12, 6149. \doi{10.1038/s41467-021-26349-x}.

\bibitem{ghgcorporate} WRI und WBCSD (2004). \emph{The Greenhouse Gas Protocol: A Corporate Accounting and Reporting Standard}, revised edition. \url{https://ghgprotocol.org/corporate-standard}.

\bibitem{ghgscope2} WRI (2015). \emph{GHG Protocol Scope 2 Guidance: An amendment to the GHG Protocol Corporate Standard}. \url{https://ghgprotocol.org/scope-2-guidance}.

\bibitem{ghgscope3} WRI und WBCSD (2011). \emph{Corporate Value Chain (Scope 3) Accounting and Reporting Standard}. \url{https://ghgprotocol.org/scope-3-standard}.

\bibitem{duchin2025} Duchin, R., Gao, J. und Xu, Q. (2025). Sustainability or Greenwashing: Evidence from the Asset Market for Industrial Pollution. \emph{The Journal of Finance}, 80(2), 699--754. Online zuerst 2024. \doi{10.1111/jofi.13412}.

\bibitem{bjorn2022} Bjørn, A., Lloyd, S. M., Brander, M. und Matthews, H. D. (2022). Renewable energy certificates threaten the integrity of corporate science-based targets. \emph{Nature Climate Change}, 12, 539--546. \doi{10.1038/s41558-022-01379-5}.

\bibitem{boulay2018} Boulay, A.-M. et al. (2018). The WULCA consensus characterization model for water scarcity footprints: assessing impacts of water consumption based on available water remaining (AWARE). \emph{The International Journal of Life Cycle Assessment}, 23, 368--378. Online zuerst 2017. \doi{10.1007/s11367-017-1333-8}.

\bibitem{rsei} U.S. EPA (o.\,J.). \emph{Understanding RSEI Results}. Offizielle Erläuterung des Risk-Screening Environmental Indicators Model. \url{https://www.epa.gov/rsei/understanding-rsei-results}. Abgerufen am 12.09.2026.

\bibitem{bls} U.S. Bureau of Labor Statistics (o.\,J.). \emph{How To Compute Your Firm's Incidence Rate for Safety Management}. \url{https://www.bls.gov/iif/overview/compute-nonfatal-incidence-rates.htm}. Recherchestand 12.09.2026.

\bibitem{delmas2011} Delmas, M. A. und Burbano, V. C. (2011). The Drivers of Greenwashing. \emph{California Management Review}, 54(1), 64--87. \doi{10.1525/cmr.2011.54.1.64}.

\bibitem{marquis2016} Marquis, C., Toffel, M. W. und Zhou, Y. (2016). Scrutiny, Norms, and Selective Disclosure: A Global Study of Greenwashing. \emph{Organization Science}, 27(2), 483--498. \doi{10.1287/orsc.2015.1039}.

\bibitem{brown2026} Brown, E., Hsu, A. und Manya, D. (2026). Red flags in green promises: a framework for identifying greenwashing risk in corporate climate pledges. \emph{npj Climate Action}, 5, 19. \doi{10.1038/s44168-026-00346-6}.

\bibitem{badgley2022} Badgley, G., Freeman, J., Hamman, J. J., Haya, B., Trugman, A. T., Anderegg, W. R. L. und Cullenward, D. (2022). Systematic over-crediting in California's forest carbon offsets program. \emph{Global Change Biology}, 28(4), 1433--1445. \doi{10.1111/gcb.15943}.

\bibitem{smaa1998} Lahdelma, R., Hokkanen, J. und Salminen, P. (1998). SMAA -- Stochastic multiobjective acceptability analysis. \emph{European Journal of Operational Research}, 106(1), 137--143. \doi{10.1016/S0377-2217(97)00163-X}.

\bibitem{constituents} DataHub / datasets (o.\,J.). \emph{S\&P 500 Companies -- Constituents}. Öffentliche Referenz-CSV. \url{https://github.com/datasets/s-and-p-500-companies}. Datenabruf 12.09.2026; Indexmitgliedschaft und Aktualität vor produktiver Verwendung prüfen.

\bibitem{sbti} Science Based Targets initiative (o.\,J.). \emph{Target Dashboard}, einschließlich By-company-Excel und Erläuterungen. \url{https://sciencebasedtargets.org/target-dashboard}. Abgerufen am 12.09.2026.

\bibitem{ca100} Climate Action 100+ (2025/2026). \emph{Net Zero Company Benchmark}. Geprüfter Excel-Export: 2025er Assessments, Version 2.2. Portal: \url{https://www.climateaction100.org/net-zero-company-benchmark/}. Abgerufen am 12.09.2026.

\bibitem{nzt2025} Tatlow, H., Manya, D. und Siu, D. (2025). \emph{Net Zero Stocktake 2025}, Version 1. Datensatz, Zenodo, veröffentlicht am 17.09.2025. \doi{10.5281/zenodo.17143240}. Lizenz: CC BY 4.0 laut Datensatzseite.

\bibitem{tpi} TPI Global Climate Transition Centre (o.\,J.). \emph{Corporate assessments: Management Quality and Carbon Performance}. \url{https://www.transitionpathwayinitiative.org/corporates}. Recherchestand 12.09.2026.

\bibitem{tpiterms} TPI Global Climate Transition Centre (o.\,J.). \emph{Use of TPI Centre data}. \url{https://transitionpathwayinitiative.org/use-of-the-centre-s-data}. Abgerufen am 12.09.2026.

\bibitem{wba} World Benchmarking Alliance (o.\,J.). \emph{Data Explorer User Guide} und \emph{API Specification}. \url{https://www.worldbenchmarkingalliance.org/documentation/user-guide}; \url{https://www.worldbenchmarkingalliance.org/documentation/api-specification}. Abgerufen am 12.09.2026.

\bibitem{epaghg} U.S. EPA (o.\,J.). \emph{Greenhouse Gas Reporting Program -- Data Sets}. \url{https://www.epa.gov/ghgreporting/data-sets}. Recherchestand 12.09.2026.

\bibitem{climatetrace} Climate TRACE (o.\,J.). \emph{Data Downloads}. \url{https://climatetrace.org/data}. Recherchestand 12.09.2026.

\bibitem{sec} U.S. Securities and Exchange Commission (o.\,J.). \emph{EDGAR Application Programming Interfaces}. \url{https://www.sec.gov/search-filings/edgar-application-programming-interfaces}. Beschreibung der CompanyFacts-Schnittstelle; direkter Seitenabruf in dieser Umgebung eingeschränkt.

\bibitem{pcaf2025} PCAF (2025). \emph{The Global GHG Accounting and Reporting Standard for the Financial Industry, Part A: Financed Emissions}, 3rd edition; einschließlich separat bereitgestellter Supplemental Guidance on Financed Avoided Emissions and Forward-Looking Metrics. \url{https://carbonaccountingfinancials.com/standard}. Abgerufen am 12.09.2026.

\bibitem{fraser2023} Fraser, A. und Fiedler, T. (2023). Net-zero targets for investment portfolios: An analysis of financed emissions metrics. \emph{Energy Economics}, 126, 106917. \doi{10.1016/j.eneco.2023.106917}.

\bibitem{koelbel2020} Kölbel, J. F., Heeb, F., Paetzold, F. und Busch, T. (2020). Can Sustainable Investing Save the World? Reviewing the Mechanisms of Investor Impact. \emph{Organization \& Environment}, 33(4), 554--574. \doi{10.1177/1086026620919202}.

\bibitem{hartzmark2026} Hartzmark, S. M. und Shue, K. (2026). \emph{Counterproductive Sustainable Investing: The Impact Elasticity of Brown and Green Firms}. NBER Working Paper 35519, Juli 2026. Nicht als Journalpublikation zitiert. \doi{10.3386/w35519}.

\bibitem{pastor2021} Pástor, Ľ., Stambaugh, R. F. und Taylor, L. A. (2021). Sustainable investing in equilibrium. \emph{Journal of Financial Economics}, 142(2), 550--571. \doi{10.1016/j.jfineco.2020.12.011}.
\end{thebibliography}
\end{document}
